% Options for packages loaded elsewhere
% Options for packages loaded elsewhere
\PassOptionsToPackage{unicode}{hyperref}
\PassOptionsToPackage{hyphens}{url}
\PassOptionsToPackage{dvipsnames,svgnames,x11names}{xcolor}
%
\documentclass[
  letterpaper,
  DIV=11,
  numbers=noendperiod]{scrartcl}
\usepackage{xcolor}
\usepackage{amsmath,amssymb}
\setcounter{secnumdepth}{5}
\usepackage{iftex}
\ifPDFTeX
  \usepackage[T1]{fontenc}
  \usepackage[utf8]{inputenc}
  \usepackage{textcomp} % provide euro and other symbols
\else % if luatex or xetex
  \usepackage{unicode-math} % this also loads fontspec
  \defaultfontfeatures{Scale=MatchLowercase}
  \defaultfontfeatures[\rmfamily]{Ligatures=TeX,Scale=1}
\fi
\usepackage{lmodern}
\ifPDFTeX\else
  % xetex/luatex font selection
\fi
% Use upquote if available, for straight quotes in verbatim environments
\IfFileExists{upquote.sty}{\usepackage{upquote}}{}
\IfFileExists{microtype.sty}{% use microtype if available
  \usepackage[]{microtype}
  \UseMicrotypeSet[protrusion]{basicmath} % disable protrusion for tt fonts
}{}
\makeatletter
\@ifundefined{KOMAClassName}{% if non-KOMA class
  \IfFileExists{parskip.sty}{%
    \usepackage{parskip}
  }{% else
    \setlength{\parindent}{0pt}
    \setlength{\parskip}{6pt plus 2pt minus 1pt}}
}{% if KOMA class
  \KOMAoptions{parskip=half}}
\makeatother
% Make \paragraph and \subparagraph free-standing
\makeatletter
\ifx\paragraph\undefined\else
  \let\oldparagraph\paragraph
  \renewcommand{\paragraph}{
    \@ifstar
      \xxxParagraphStar
      \xxxParagraphNoStar
  }
  \newcommand{\xxxParagraphStar}[1]{\oldparagraph*{#1}\mbox{}}
  \newcommand{\xxxParagraphNoStar}[1]{\oldparagraph{#1}\mbox{}}
\fi
\ifx\subparagraph\undefined\else
  \let\oldsubparagraph\subparagraph
  \renewcommand{\subparagraph}{
    \@ifstar
      \xxxSubParagraphStar
      \xxxSubParagraphNoStar
  }
  \newcommand{\xxxSubParagraphStar}[1]{\oldsubparagraph*{#1}\mbox{}}
  \newcommand{\xxxSubParagraphNoStar}[1]{\oldsubparagraph{#1}\mbox{}}
\fi
\makeatother


\usepackage{longtable,booktabs,array}
\newcounter{none} % for unnumbered tables
\usepackage{calc} % for calculating minipage widths
% Correct order of tables after \paragraph or \subparagraph
\usepackage{etoolbox}
\makeatletter
\patchcmd\longtable{\par}{\if@noskipsec\mbox{}\fi\par}{}{}
\makeatother
% Allow footnotes in longtable head/foot
\IfFileExists{footnotehyper.sty}{\usepackage{footnotehyper}}{\usepackage{footnote}}
\makesavenoteenv{longtable}
\usepackage{graphicx}
\makeatletter
\newsavebox\pandoc@box
\newcommand*\pandocbounded[1]{% scales image to fit in text height/width
  \sbox\pandoc@box{#1}%
  \Gscale@div\@tempa{\textheight}{\dimexpr\ht\pandoc@box+\dp\pandoc@box\relax}%
  \Gscale@div\@tempb{\linewidth}{\wd\pandoc@box}%
  \ifdim\@tempb\p@<\@tempa\p@\let\@tempa\@tempb\fi% select the smaller of both
  \ifdim\@tempa\p@<\p@\scalebox{\@tempa}{\usebox\pandoc@box}%
  \else\usebox{\pandoc@box}%
  \fi%
}
% Set default figure placement to htbp
\def\fps@figure{htbp}
\makeatother





\setlength{\emergencystretch}{3em} % prevent overfull lines

\providecommand{\tightlist}{%
  \setlength{\itemsep}{0pt}\setlength{\parskip}{0pt}}



 
\usepackage[backend=biber]{biblatex}
\addbibresource{AIM-143-1999-284.bib}


\KOMAoption{captions}{tableheading}
\makeatletter
\@ifpackageloaded{tcolorbox}{}{\usepackage[skins,breakable]{tcolorbox}}
\@ifpackageloaded{fontawesome5}{}{\usepackage{fontawesome5}}
\definecolor{quarto-callout-color}{HTML}{909090}
\definecolor{quarto-callout-note-color}{HTML}{0758E5}
\definecolor{quarto-callout-important-color}{HTML}{CC1914}
\definecolor{quarto-callout-warning-color}{HTML}{EB9113}
\definecolor{quarto-callout-tip-color}{HTML}{00A047}
\definecolor{quarto-callout-caution-color}{HTML}{FC5300}
\definecolor{quarto-callout-color-frame}{HTML}{acacac}
\definecolor{quarto-callout-note-color-frame}{HTML}{4582ec}
\definecolor{quarto-callout-important-color-frame}{HTML}{d9534f}
\definecolor{quarto-callout-warning-color-frame}{HTML}{f0ad4e}
\definecolor{quarto-callout-tip-color-frame}{HTML}{02b875}
\definecolor{quarto-callout-caution-color-frame}{HTML}{fd7e14}
\makeatother
\makeatletter
\@ifpackageloaded{caption}{}{\usepackage{caption}}
\AtBeginDocument{%
\ifdefined\contentsname
  \renewcommand*\contentsname{Table of contents}
\else
  \newcommand\contentsname{Table of contents}
\fi
\ifdefined\listfigurename
  \renewcommand*\listfigurename{List of Figures}
\else
  \newcommand\listfigurename{List of Figures}
\fi
\ifdefined\listtablename
  \renewcommand*\listtablename{List of Tables}
\else
  \newcommand\listtablename{List of Tables}
\fi
\ifdefined\figurename
  \renewcommand*\figurename{Figure}
\else
  \newcommand\figurename{Figure}
\fi
\ifdefined\tablename
  \renewcommand*\tablename{Table}
\else
  \newcommand\tablename{Table}
\fi
}
\@ifpackageloaded{float}{}{\usepackage{float}}
\floatstyle{ruled}
\@ifundefined{c@chapter}{\newfloat{codelisting}{h}{lop}}{\newfloat{codelisting}{h}{lop}[chapter]}
\floatname{codelisting}{Listing}
\newcommand*\listoflistings{\listof{codelisting}{List of Listings}}
\usepackage{amsthm}
\theoremstyle{remark}
\AtBeginDocument{\renewcommand*{\proofname}{Proof}}
\newtheorem*{remark}{Remark}
\newtheorem*{solution}{Solution}
\newtheorem{refremark}{Remark}[section]
\newtheorem{refsolution}{Solution}[section]
\makeatother
\makeatletter
\makeatother
\makeatletter
\@ifpackageloaded{caption}{}{\usepackage{caption}}
\@ifpackageloaded{subcaption}{}{\usepackage{subcaption}}
\makeatother
\usepackage{bookmark}
\IfFileExists{xurl.sty}{\usepackage{xurl}}{} % add URL line breaks if available
\urlstyle{same}
% fallback for those not using the hyperref driver hyperxmp:
\makeatletter
\@ifundefined{xmpquote}{\newcommand{\xmpquote}[1]{#1}}{}
\makeatother
\hypersetup{
  pdftitle={Notes on the Riemann zeta function{,} 2},
  pdfauthor={Michel Balazard, Eric Saias, and Marc Yor},
  colorlinks=true,
  linkcolor={blue},
  filecolor={Maroon},
  citecolor={Blue},
  urlcolor={Blue},
  pdfcreator={LaTeX via pandoc}}


\title{Notes on the Riemann zeta function, 2}
\author{Michel Balazard, Eric Saias, and Marc Yor}
\date{1999}
\begin{document}
\maketitle

\renewcommand*\contentsname{Table of contents}
{
\setcounter{tocdepth}{3}
\tableofcontents
}

\emph{This document is a translation into English of the following:}

\begin{quote}
Balazard, M., Saias, E., and Yor, M. ``Notes sur la fonction \(\zeta\)
de Riemann, 2''. \emph{Advances in Mathematics} \textbf{143} (1999),
284--287.
\end{quote}

\emph{The translator (\href{https://thosgood.net}{Tim Hosgood}) takes
full responsibility for any errors introduced in this document, and
claims no rights to any of the mathematical content.}

\tableofcontents

\providecommand{\scr}[1]{{\mathscr{#1}}}
\renewcommand{\cal}[1]{{\mathcal{#1}}}
\renewcommand{\frak}[1]{{\mathfrak{#1}}}
\renewcommand{\geq}{\geqslant}
\renewcommand{\leq}{\leqslant}

\newcommand{\dd}{\operatorname{d}\!}

\section*{Introduction}\label{introduction}
\addcontentsline{toc}{section}{Introduction}

\oldpage{284}

We denote by \(\sum_{\Re\rho>1/2}\) a sum over the possible zeros of
\(\zeta(s)\) with real part greater than \(\frac12\), where the zeros of
multiplicity \(m\) are counted \(m\) times. The goal of this note is the
proof of the following result.

\begin{itenv}{Theorem}
We have \[
  \frac{1}{2\pi}\int_{\Re(s)=1/2} \frac{\log|\zeta(s)|}{|s|^2}|\operatorname{d}\!s|
  = \sum_{\Re\rho>1/2} \log\left|\frac{\rho}{1-\rho}\right|.
\tag{1}
\] In particular, the Riemann hypothesis is true if and only if \[
  \int_{\Re(s)=1/2} \frac{\log|\zeta(s)|}{|s|^2}|\operatorname{d}\!s| = 0.
\]

\end{itenv}

\begin{proof}
This proof consists of two steps.

\emph{First step.} We start by stating some properties satisfied by a
generic function \(f\) in the Hardy space \(H^p(\mathbf{D})\), where
\(\mathbf{D}=\{z\in\mathbb{C}:|z|<1\}\), and \(p\) is a positive real
number. We denote by \(f^*\) the function defined almost everywhere on
the trigonometric circle \(\partial\mathbf{D}=\{z\in\mathbb{C}:|z|=1\}\)
by \(f^*(e^{i\theta})=\lim_{r\to1^-}f(re^{i\theta})\). We use the letter
\(z\) to denote an element of the trigonometric disc \(\mathbf{D}\), and
write \[
  s = s(z) = \frac12+\frac{1+z}{2(1-z)} = \frac{1}{1-z}.
\] This formula defines a conformal representation of the disc
\(\mathbf{D}\) in the semi-plane \(\Re(s)>1/2\).

\oldpage{285}

By Jensen's formula (see, for example, \autocite[Theorem 3.61]{4}), we
have, for \(f(0)\neq0\) and \(r<1\), \[
  \frac{1}{2\pi}\int_{-\pi}^\pi \log|f(re^{i\theta})|\operatorname{d}\!\theta
  = \log|f(0)| + \sum_{\substack{|\alpha|<r\\f(\alpha)=0}} \log\frac{r}{|a|}
\tag{2}
\] where, in the sum, the zeros of multiplicity \(m\) are counted \(m\)
times. Denote by \[
  \exp\left\{
    -\int_{-\pi}^\pi \frac{e^{i\theta}+z}{e^{i\theta}-z}\operatorname{d}\!\mu(\theta)
  \right\}
\] the singular interior factor of \(f\). As \(r\) tends to \(1\),
Equation (2) becomes (cf. \autocite[68]{2}) \[
  \frac{1}{2\pi}\int_{-\pi}^\pi \log|f(re^{i\theta})|\operatorname{d}\!\theta
  = \log|f(0)| + \sum_{\substack{|\alpha|<1\\f(\alpha)=0}} \log\frac{1}{|a|} + \int_{-\pi}^\pi\operatorname{d}\!\mu(\theta).
\tag{3}
\] This formula is a consequence of the factorisation theorem for
functions in \(H^p\); it is stated in \autocite{2} for \(p=1\), but also
holds for all positive values of \(p\).

\emph{Second step.} Now consider the function \[
  f(z) = (s-1)\zeta(s)
\] (where \(s=1/(1-z)\)). The elementary properties of the Riemann
\(\zeta\) function (see, for example, \autocite{5}) allow us to show
that, on one hand, \(f\) belongs to the Hardy space
\(H^{1/3}(\mathbf{D})\), and, on the other hand, that the measure
\(\mu\) associated to the singular interior factor of \(f\) is zero (for
this latter point, it suffices to reuse the argument developed by
Bercovici and Foias for the interior factor of the functions
\((\theta-\theta^s)\zeta(s)(s+1/2)/s\), found in the proof of
\autocite[Proposition 2.1]{1}). We can equally show that \[
  \begin{aligned}
    \int_{-\pi}^\pi \log|f^*(e^{i\theta})|\operatorname{d}\!\theta
    &= \int_{\Re(s)=1/2} \frac{\log|\zeta(s)|}{|s|^2}|\operatorname{d}\!s|,
  \\\log|f(0)|
    &= 0,
  \\\sum_{\substack{|\alpha|<1\\f(\alpha)=0}} \log\frac{1}{|\alpha|}
    &= \sum_{\Re\rho>1/2} \log\left|\frac{\rho}{1-\rho}\right|.
  \end{aligned}
\] With all this information, our result follows from Equation (3).
\end{proof}

\oldpage{286}

We finish with some remarks. There are statements related to ours in the
works \autocite{7,6} of Wang and Volchkov. It is even possible that
Jensen himself was aware of Equation (1) (the reader can consult the
article \autocite{3} where Jensen informs Mittag--Leffler of his
discovery of Equation (2)). It seem interesting, however, to present
things as we have done here, and this is for the following three
reasons:

\begin{enumerate}
\def\labelenumi{\alph{enumi}.}
\tightlist
\item
  Equation (1) is simpler than those that appear in \autocite{7,6};
\item
  we show here that, to establish Equation (1), it is natural to place
  ourselves in the framework of Hardy spaces;
\item
  the form of the integral in Equation (1) allows us to interpret this
  result via Brownian motion, as we show below.
\end{enumerate}

Denote by \(Z=X+iY\) the planar Brownian motion from \(0\) (or from
\(1\)), and by \(Z_{T_{1/2}}=\frac12+iY_{T_{1/2}}\) its first point of
impact on the critical line \(\Re s=1/2\), where
\(T_{1/2}:=\inf\{t:X_t=1/2\}\). We know that \(Y_{T_{1/2}}\) follows a
Cauchy law with parameter \(1/2\). In other words, the law of
\(Y_{T_{1/2}}\) has density \(1/2\pi(1/4+t^2)\). Thus the second part of
the theorem can be stated in the following manner: the Riemann
hypothesis is true if and only if \[
  \mathbb{E}[\log\vert\zeta(Z_{T_{1/2}})] = 0.
\]

\subsection*{Thanks}\label{thanks}
\addcontentsline{toc}{subsection}{Thanks}

We thank Luis Báez-Duarte, Michel Delasneri, Catherine Donati, Laurent
Habsieger, Aleksandar Ivić, and Alain Plagne for useful conversations.


\nocite{*}
\printbibliography



\end{document}
